\documentclass[a4paper]{article}
\usepackage{amsfonts}
\usepackage{amsmath}
\usepackage{amssymb}
\usepackage{graphicx}     

\newcommand{\dint}{\displaystyle\int}

\begin{document}
  
\noindent \large Auteur: Thijs Van Schel \\
\noindent \large Studierichting: Informatica\\
\noindent \large Oefeningenreeks 4A nr. 58.12\\

\medskip

\normalsize

$\dint \dfrac{dx}{3 \sin x + 2 \cos x + 2}$ \\

$t = \tan \dfrac{x}{2} \Rightarrow \sin x = \dfrac{2t}{1+t^2}, \cos x = \dfrac{1-t^2}{1+t^2}, dx = \dfrac{2 dt}{1+t^2}$ \\

$= \dint \dfrac{\dfrac{2}{1+t^2} dt}{3 \left( \dfrac{2t}{1+t^2} \right) + 2 \left( \dfrac{1-t^2}{1+t^2} \right) + 2}$ \\

$= \dint \dfrac{2 dt}{6t + 2(1-t^2) + 2(1+t^2)}$ \\

$= \dint \dfrac{2 dt}{6t + 2 - 2t^2 + 2 + 2t^2}$ \\

$= \dint \dfrac{2 dt}{6t + 4}$ \\

$= \dint \dfrac{dt}{3t + 2}$ \\

$= \dfrac{1}{3} \ln |3t + 2| + C$ \\

$= \dfrac{1}{3} \ln \left| 3 \tan \dfrac{x}{2} + 2 \right| + C$

\begin{thebibliography}{999}
%\bibitem{a} Referentie
\end{thebibliography}
\end{document} 